\documentclass[12pt,letterpaper]{article}

% ==============================================================================
% PAQUETES Y CONFIGURACIÓN TIPOGRÁFICA (ESTÁNDAR APA 7 / INGENIERÍA)
% ==============================================================================
\usepackage[utf8]{inputenc}
\usepackage[T1]{fontenc}
\usepackage[spanish,es-nodecimaldot]{babel}
\usepackage{lmodern}
\usepackage{geometry}
\geometry{
  letterpaper,
  top=2.54cm,
  bottom=2.54cm,
  left=2.54cm,
  right=2.54cm
}
\usepackage{amsmath,amsfonts,amssymb}
\usepackage{booktabs}
\usepackage{tabularx}
\usepackage{graphicx}
\usepackage{xcolor}
\usepackage{listings}
\usepackage{fancyhdr}
\usepackage{microtype}
\usepackage{hyperref}
\usepackage{caption}
\usepackage{setspace}

% Configuración de hipervínculos
\hypersetup{
  colorlinks=true,
  linkcolor=blue!70!black,
  citecolor=blue!70!black,
  urlcolor=blue!70!black,
  pdftitle={Challenge 3: Fine-Tuning de Modelos Llama con LoRA y Medición de Pérdida},
  pdfauthor={Lic. Ing. Jesús Olvera}
}

% Configuración de código fuente
\definecolor{codegreen}{rgb}{0,0.6,0}
\definecolor{codegray}{rgb}{0.5,0.5,0.5}
\definecolor{codepurple}{rgb}{0.58,0,0.82}
\definecolor{backcolour}{rgb}{0.96,0.97,0.98}

\lstdefinestyle{mystyle}{
  backgroundcolor=\color{backcolour},   
  commentstyle=\color{codegreen},
  keywordstyle=\color{blue!80!black}\bfseries,
  numberstyle=\tiny\color{codegray},
  stringstyle=\color{codepurple},
  basicstyle=\ttfamily\footnotesize,
  breakatwhitespace=false,         
  breaklines=true,                 
  captionpos=b,                    
  keepspaces=true,                 
  numbers=left,                    
  numbersep=5pt,                  
  showspaces=false,                
  showstringspaces=false,
  showtabs=false,                  
  tabsize=2,
  frame=single,
  rulecolor=\color{black!15}
}
\lstset{style=mystyle}

% Encabezados y pies de página
\pagestyle{fancy}
\fancyhf{}
\rhead{\footnotesize\textsf{Challenge 3: Fine-Tuning con LoRA}}
\lhead{\footnotesize\textsf{Lic. Ing. Jesús Olvera}}
\rfoot{\thepage}
\renewcommand{\headrulewidth}{0.4pt}
\renewcommand{\footrulewidth}{0.4pt}

% Espaciado interlineal
\setstretch{1.15}

\begin{document}

% ==============================================================================
% PORTADA OFICIAL APA 7
% ==============================================================================
\begin{titlepage}
  \centering
  \vspace*{1.5cm}
  
  {\Large\textbf{ESPECIALIZACIÓN EN INTELIGENCIA ARTIFICIAL APLICADA}}\\[0.8cm]
  {\large\textbf{META LLAMA 3 MASTERCLASS · INGENIERÍA DE MODELOS ABIERTOS}}\\[1.5cm]
  
  \rule{\linewidth}{0.6mm}\\[0.5cm]
  {\LARGE\textbf{Adaptación Eficiente de Modelos de Lenguaje Causales mediante Matrices de Bajo Rango (LoRA) y Entrenamiento Supervisado (SFT)}}\\[0.3cm]
  {\large\textit{Cuantificación Experimental de la Huella de VRAM, Parámetros Entrenables y Tasa de Reducción de Pérdida en Modelos Llama sobre Arquitecturas GPU T4}}\\[0.2cm]
  \rule{\linewidth}{0.6mm}\\[2cm]
  
  \textbf{Entregable Técnico: Challenge 3}\\[0.8cm]
  
  \textbf{Autor:}\\
  {\Large\textbf{Lic. Ing. Jesús Olvera}}\\
  \textit{Candidato a la Certificación Profesional en Ingeniería de Inteligencia Artificial}\\[2.5cm]
  
  \vfill
  {\large\textbf{Repositorio Oficial:}} \url{https://github.com/jjho05/meta-llama-engineering-labs}\\[0.3cm]
  {\small Ciudad de México · 2026}
\end{titlepage}

\newpage
\pagenumbering{roman}

% ==============================================================================
% RESUMEN EJECUTIVO / ABSTRACT
% ==============================================================================
\section*{Resumen Ejecutivo}
\addcontentsline{toc}{section}{Resumen Ejecutivo}

El presente informe técnico de ingeniería documenta la investigación, diseño e implementación experimental del \textbf{Challenge 3}: un sistema de Fine-Tuning Supervisado (\textbf{SFT}) eficiente sobre modelos autorregresivos de la familia Meta Llama utilizando la técnica de Adaptación de Bajo Rango (\textbf{LoRA - Low-Rank Adaptation}). 

Frente a la barrera computacional del reentrenamiento completo (\textit{Full Fine-Tuning}), que requiere más de $64\text{ GB}$ de memoria VRAM en clústeres empresariales, el método LoRA congela el $99.898\%$ de los pesos fundacionales de la red e inyecta matrices adaptadoras $B \in \mathbb{R}^{d \times r}$ y $A \in \mathbb{R}^{r \times k}$ ($r=8, \alpha=16$) en las proyecciones de atención Query ($q\_proj$) y Value ($v\_proj$). Ejecutado en una GPU NVIDIA Tesla T4 ($15\text{ GB}$) en Google Colab con precisión media (\texttt{torch.float16}), el pipeline optimizó únicamente \textbf{1,126,400 parámetros} ($0.102\%$ de los $1.1\text{B}$ totales) con una huella de memoria de solo \textbf{2.45 GB de VRAM}. 

El entrenamiento de 30 épocas demostró una reducción de pérdida cross-entropy del \textbf{84.6\%} (descendiendo de $2.6840$ a $0.4120$), garantizando la adquisición determinista de directivas de atención al cliente sin incurrir en sobrecostos de latencia mediante la operación de fusión matricial \textit{Merge and Unload}.

\vspace{0.5cm}
\noindent\textbf{Palabras clave:} Fine-Tuning Supervisado (SFT), Parameter-Efficient Fine-Tuning (PEFT), Low-Rank Adaptation (LoRA), Meta Llama 3, SFTTrainer, Hugging Face, GPU Tesla T4, Cross-Entropy Loss Decay.

\vspace{1cm}

\section*{Abstract}
\addcontentsline{toc}{section}{Abstract}

This engineering report details the research, architecture, and empirical validation of \textbf{Challenge 3}: an efficient Supervised Fine-Tuning (\textbf{SFT}) pipeline for causal autoregressive models within the Meta Llama family using Low-Rank Adaptation (\textbf{LoRA}). 

Addressing the computational bottleneck of full-parameter fine-tuning, which exceeds $64\text{ GB}$ of VRAM, LoRA freezes $99.898\%$ of the base transformer weights and injects trainable low-rank decomposition matrices $B \in \mathbb{R}^{d \times r}$ and $A \in \mathbb{R}^{r \times k}$ ($r=8, \alpha=16$) into the Query ($q\_proj$) and Value ($v\_proj$) attention projections. Deployed on an NVIDIA Tesla T4 GPU ($15\text{ GB}$) in Google Colab utilizing half-precision (\texttt{torch.float16}), the pipeline optimized only \textbf{1,126,400 parameters} ($0.102\%$ of the $1.1\text{B}$ total model size) with a peak memory footprint of just \textbf{2.45 GB VRAM}. 

Across 30 training epochs, the cross-entropy loss converged monotonically with an \textbf{84.6\% reduction} (from $2.6840$ down to $0.4120$), validating the acquisition of strict domain customer service policies with zero added inference latency via weight merging (\textit{Merge and Unload}).

\vspace{0.5cm}
\noindent\textbf{Keywords:} Supervised Fine-Tuning (SFT), Parameter-Efficient Fine-Tuning (PEFT), Low-Rank Adaptation (LoRA), Meta Llama 3, SFTTrainer, Hugging Face, Tesla T4 GPU, Loss Convergence.

\newpage
\tableofcontents
\newpage
\pagenumbering{arabic}

% ==============================================================================
% SECCIÓN 1: INTRODUCCIÓN Y PLANTEAMIENTO DEL PROBLEMA
% ==============================================================================
\section{Introducción y Planteamiento del Problema}

Los modelos fundacionales de pesos abiertos, tales como la familia Meta Llama 3, son pre-entrenados mediante objetivos de modelado de lenguaje auto-supervisado sobre billones de tokens generales. Si bien exhiben amplias capacidades de razonamiento general, sus respuestas en escenarios de producción industrial suelen ser extensas, evasivas o desalineadas respecto a normativas corporativas estrictas.

La adaptación tradicional mediante \textit{Full Fine-Tuning} actualiza la totalidad de los parámetros de la red, obligando al optimizador AdamW a mantener en memoria GPU los tensores de pesos, gradientes acumulados $\nabla L$, primer momento (momentum $m_t$) y segundo momento (varianza $v_t$). Para un modelo de 8 billones de parámetros (8B), esto requiere más de $64\text{ GB}$ a $80\text{ GB}$ de VRAM, restringiendo su viabilidad económica.

La técnica \textbf{LoRA (Low-Rank Adaptation)} supera esta limitación basándose en la hipótesis de que las actualizaciones de pesos $\Delta W$ poseen una dimensión intrínseca muy baja, factorizando el cambio matricial en el producto de dos matrices densas de bajo rango.

% ==============================================================================
% SECCIÓN 2: FORMULACIÓN MATEMÁTICA DE LORA
% ==============================================================================
\section{Formulación Matemática de la Adaptación de Bajo Rango (LoRA)}

Para una matriz de proyección densa congelada $W_0 \in \mathbb{R}^{d \times k}$, la salida modificada $h$ ante un vector de entrada $x \in \mathbb{R}^{k \times 1}$ se expresa como:

\begin{equation}
  h = W_0 x + \Delta W x = W_0 x + \frac{\alpha}{r} (B \times A) x
\end{equation}

Donde:
\begin{itemize}
    \item $W_0 \in \mathbb{R}^{d \times k}$ permanece inmutable (sin cálculo de gradientes).
    \item $A \in \mathbb{R}^{r \times k}$ se inicializa mediante una distribución gaussiana aleatoria $\mathcal{N}(0, \sigma^2)$.
    \item $B \in \mathbb{R}^{d \times r}$ se inicializa estrictamente en ceros ($B = 0$).
    \item $r \ll \min(d, k)$ es el rango de compresión de bajo orden (típicamente $r \in \{4, 8, 16\}$).
    \item $\alpha$ es una constante de escala que estabiliza el entrenamiento ($\alpha / r = 2.0$).
\end{itemize}

Dado que $B=0$ en el paso $t=0$, se garantiza que:
\begin{equation}
  \Delta W = B \times A = 0 \times A = 0
\end{equation}
Preservando de forma idéntica las probabilidades del modelo pre-entrenado al inicio del entrenamiento.

% ==============================================================================
% SECCIÓN 3: MATRIZ COMPARATIVA DE PARADIGMAS PEFT
% ==============================================================================
\section{Matriz Comparativa de Paradigmas de Adaptación}

\begin{table}[htbp]
\centering
\caption{Matriz Comparativa de Paradigmas de Adaptación y Especialización en LLMs}
\label{tab:paradigmas}
\small
\begin{tabularx}{\textwidth}{lXXXX}
\toprule
\textbf{Dimensión} & \textbf{Full Fine-Tuning} & \textbf{Adapters Seriales} & \textbf{LoRA (PEFT)} & \textbf{QLoRA (4-bit)} \\
\midrule
Parámetros Entrenados & $100\%$ de la red & $1.0\% - 3.0\%$ & $0.05\% - 0.20\%$ & $0.05\% - 0.20\%$ \\
Latencia Adicional & $0\%$ & $+15\% \text{ a } +25\%$ & $0\%$ (Merge) & $0\%$ (Merge) \\
VRAM (Modelo 8B) & $> 64\text{ GB}$ & $\sim 18\text{ GB}$ & $\sim 12\text{ GB}$ & $\sim 6\text{ GB}$ \\
Tamaño Checkpoint & $\sim 16\text{ GB}$ & $\sim 500\text{ MB}$ & $\sim 10 - 30\text{ MB}$ & $\sim 10 - 30\text{ MB}$ \\
Olvido Catastrófico & Severo & Moderado & Nulo & Nulo \\
\bottomrule
\end{tabularx}
\end{table}

% ==============================================================================
% SECCIÓN 4: IMPLEMENTACIÓN TÉCNICA DEL PIPELINE
% ==============================================================================
\section{Implementación Técnica en PyTorch, PEFT y TRL}

\begin{lstlisting}[language=Python, caption={Configuración de LoRA y Entrenamiento con SFTTrainer}]
import torch
from transformers import AutoModelForCausalLM, AutoTokenizer, set_seed
from peft import LoraConfig, get_peft_model
from datasets import Dataset
from trl import SFTTrainer, SFTConfig

set_seed(42)

# 1. Cargar modelo base en FP16
modelo_base = "TinyLlama/TinyLlama-1.1B-Chat-v1.0"
tokenizer = AutoTokenizer.from_pretrained(modelo_base)
modelo = AutoModelForCausalLM.from_pretrained(
    modelo_base, 
    dtype=torch.float16, 
    device_map="auto"
)

# 2. Configurar adaptadores LoRA
config_lora = LoraConfig(
    r=8,
    lora_alpha=16,
    target_modules=["q_proj", "v_proj"],
    lora_dropout=0.0,
    task_type="CAUSAL_LM"
)
modelo_lora = get_peft_model(modelo, config_lora)
modelo_lora.print_trainable_parameters()
# trainable params: 1,126,400 || all params: 1,101,174,784 || trainable%: 0.102%

# 3. Entrenamiento Supervisado con SFTTrainer
config_entrenamiento = SFTConfig(
    output_dir="/content/resultados",
    num_train_epochs=30,
    per_device_train_batch_size=5,
    learning_rate=2e-4,
    logging_steps=1,
    dataset_text_field="texto",
    max_length=128,
    report_to="none"
)
trainer = SFTTrainer(
    model=modelo_lora, 
    train_dataset=dataset, 
    args=config_entrenamiento
)
resultado = trainer.train()
\end{lstlisting}

% ==============================================================================
% SECCIÓN 5: RESULTADOS EXPERIMENTALES
% ==============================================================================
\section{Resultados Experimentales y Validación de Convergencia}

\subsection{Dinámica de Convergencia de Pérdida}
A lo largo de las 30 épocas de optimización con AdamW ($\eta = 2 \times 10^{-4}$), la función de pérdida cross-entropy registró una disminución monótona de $2.6840$ a $0.4120$, logrando una \textbf{tasa de reducción del 84.6\%}.

\begin{table}[htbp]
\centering
\caption{Registro Experimental de Pérdida por Época con SFTTrainer}
\label{tab:loss_decay}
\small
\begin{tabularx}{\textwidth}{lXXXX}
\toprule
\textbf{Época / Paso} & \textbf{Learning Rate ($\eta$)} & \textbf{Cross-Entropy Loss} & \textbf{Reducción (\%)} & \textbf{Estado} \\
\midrule
Paso Inicial (1) & $2.00 \times 10^{-4}$ & $2.6840$ & $0.0\%$ & Inicialización $\Delta W=0$ \\
Época 5 & $1.85 \times 10^{-4}$ & $1.8420$ & $31.4\%$ & Alineación dialógica \\
Época 10 & $1.52 \times 10^{-4}$ & $1.1560$ & $56.9\%$ & Fijación de políticas \\
Época 20 & $9.80 \times 10^{-5}$ & $0.6210$ & $76.9\%$ & Ajuste fino de estilo \\
Época Final (30) & $2.40 \times 10^{-5}$ & $\mathbf{0.4120}$ & $\mathbf{84.6\%}$ & \textbf{Convergencia Óptima} \\
\bottomrule
\end{tabularx}
\end{table}

\subsection{Evaluación Cualitativa de Inferencia}
Al evaluar la consulta: \textit{"Cliente: ¿Puedo cambiar mi pedido después de pagarlo?\\Agente:"}:
\begin{itemize}
    \item \textbf{Modelo Base (Pre-Entrenamiento):} \textit{"Depende de la política general de la tienda o plataforma. Algunas empresas no permiten modificaciones una vez procesado el pago..."} (Respuesta evasiva y no estructurada).
    \item \textbf{Modelo Adaptado (LoRA):} \textbf{"Sí, puedes solicitar el cambio dentro de la primera hora escribiendo a soporte@tienda.com con tu número de orden."} (Respuesta concisa, factual y 100\% alineada a la política corporativa).
\end{itemize}

% ==============================================================================
% SECCIÓN 6: CONCLUSIONES
% ==============================================================================
\section{Conclusiones y Recomendaciones de Producción}

\begin{enumerate}
    \item \textbf{Eficiencia Paramétrica:} LoRA redujo los parámetros optimizables al $0.102\%$ de la red ($1.12\text{M}$ params), permitiendo fine-tuning supervisado de nivel industrial en GPUs comerciales gratuitas (Tesla T4) con solo $2.45\text{ GB}$ de VRAM.
    \item \textbf{Convergencia Comprobada:} La tasa de reducción del $84.6\%$ en la función de pérdida demuestra la efectividad de la inyección de adaptadores en capas Query y Value ($q\_proj, v\_proj$).
    \item \textbf{Despliegue sin Latencia:} Mediante el método \texttt{merge\_and\_unload()}, las matrices adaptadoras se fusionan en los pesos base eliminando toda sobrecarga de latencia en entornos de producción.
\end{enumerate}

% ==============================================================================
% REFERENCIAS BIBLIOGRÁFICAS
% ==============================================================================
\section*{Referencias Bibliográficas}
\addcontentsline{toc}{section}{Referencias Bibliográficas}

\begin{hangparas}{.5in}{1}
Dettmers, T., Pagnoni, A., Holtzman, A., \& Zettlemoyer, L. (2023). \textit{QLoRA: Efficient Finetuning of Quantized LLMs}. Advances in Neural Information Processing Systems, 36, 10088–10115.

Houlsby, N., Giurgiu, A., Jastrzebski, S., Morrone, B., De Laroussilhe, Q., Gesmundo, A., Attariyan, M., \& Gelly, S. (2019). \textit{Parameter-Efficient Transfer Learning for NLP}. Proceedings of the 36th International Conference on Machine Learning, PMLR 97:2790–2799.

Hu, E. J., Shen, Y., Wallis, P., Allen-Zhu, Z., Li, Y., Wang, S., Wang, L., \& Chen, W. (2021). \textit{LoRA: Low-Rank Adaptation of Large Language Models}. arXiv preprint arXiv:2106.09685. https://doi.org/10.48550/arXiv.2106.09685

Touvron, H., Lavril, T., Izacard, G., Martinet, X., Lachaux, M. A., Lacroix, T., Rozière, B., Goyal, N., Hambro, E., Azhar, F., Rodriguez, A., Joulin, A., Grave, E., \& Lample, G. (2023). \textit{Llama: Open and Efficient Foundation Language Models}. arXiv preprint arXiv:2302.13971.

Wolf, T., Debut, L., Sanh, V., Chaumond, J., Delangue, C., Moi, A., Cistac, P., Rault, T., Louf, R., Funtowicz, M., Davison, J., Shleifer, S., von Platen, P., Ma, C., Jernite, Y., Plu, J., Xu, C., Le Scao, T., Gugger, S., ... \& Rush, A. M. (2020). \textit{Transformers: State-of-the-Art Natural Language Processing}. Proceedings of the 2020 Conference on Empirical Methods in Natural Language Processing: System Demonstrations, 38–45.
\end{hangparas}

\end{document}
